\MajorSection{القسم الثالث}

\Couplet{٣٩}{أُفٍّ لكم، يا أشياء الخيال، يا ذرات الغبار الراقصة في وهج الشمس؛}{يا من تؤسسون الأزليات وتبنونها على أقصر لحظة هنا في الأسفل.}{}

\Couplet{٤٠}{يا من تعبرون الحياة كطيور في أقفاص، أسرى لإرادة مستبدة؛}{ولا تفتؤون تتساءلون: كيف، ومتى، ولماذا، ومن أين، وإلى أين؟ ولا تفتؤون.}{}

\Couplet{٤١}{ولا تفتؤون تتساءلون كيف جاءت هذه الأعجوبة، لأن حيوانين من الثدييات تزاوجا واختارا}{أن يرويا عطش الحب الجسدي، فقام بذلك «الكائن الخالد».}{}

\Couplet{٤٢}{تتأملون في عجب الرضيع ذا العينين المحدّقتين، المُخرَج قسرًا من الليل إلى النهار؛}{وقد أمسكت به قبضة الحياة العملاقة، كغبار ولدته العاصفة أو رذاذ اعتصرته الريح.}{}

\Couplet{٤٣}{يأتي العالم عاجزًا، وسط خطر مزدوج وأنين ودموع؛}{لعبةً وأضحوكةً، لقيطًا وتائهًا في أيدي الأهواء والخطأ والغضب والمخاوف.}{}

\Couplet{٤٤}{لا يعرف من أين جاء ولا لماذا، ولا يدري إلى أين يتجه ولا متى؛}{ومع ذلك، فهذه هي أصفى عطايا الله، والبركة التي يحلم بها الحمقى.}{}

\Couplet{٤٥}{ثم يعود خطوة خطوة، قسرًا، إلى طفولة لا دراية لها، شاحبًا أبيض باردًا؛}{يلثغ من جديد بكلماته المتكسرة، إلى أن تُروى الحكاية كلها.}{}

\Couplet{٤٦}{وتتأملون في عجب الرضيعَ الذي انطفأت عيناه، شيخًا حنته السنوات الثقيلة؛}{كيف نجت الزورقة من مئة عاصفة، وكيف نجا الخيط من ألف مقص.}{}

\Couplet{٤٧}{وكيف جاء إلى المأدبة بغير دعوة، فوجد المائدة الفخمة مبسوطة}{بثمار سدوم الحسنة المظهر، وبحجارة تتخذ هيئة الخبز.}{}

\Couplet{٤٨}{وكيف لم تكن الحياة إلا شعاع شمس شق ظلمة كثيفة عمياء؛}{إلا هذيان العاصفة الطائشة، وصراخ الريح النهمة.}{}

\Couplet{٤٩}{وكيف خدعت نومه رؤى جميلة، لا تنفك تتلاشى عند انبثاق الصباح؛}{إلى أن صار كل حلو مرًّا، وصارت كل وردة شوكة.}{}

\Couplet{٥٠}{إلى أن لقيت عيناه التراب والرماد حيثما استدار نظرهما المحزون؛}{وحطام المسرّات والآمال والأحبّة، وركام أيامه المهدورة.}{}

\Couplet{٥١}{وكيف أخفق كل فكر بطولي سامٍ اشتهى أن يتنفس هواء السماء العليا؛}{فلم تنبت له ريشاته، وسقط إلى الأرض، وهلك من اليأس المحض.}{}

\Couplet{٥٢}{وكيف، وقد أُوتي ميراث دماغ بلغت قوته أن حلّل الشعاع الشمسي؛}{فإن بقيّته أغلظ التراب وأخشنه، كتاج من ذهب على جبين من طين.}{}

\Couplet{٥٣}{هذا البيت، هيكله لحم وعظم، وملاطه دم، وواجهته جلد؛}{موطن المرض والأوجاع والشيخوخة، قذر من الخارج، مدنّس من الداخل.}{}

\Couplet{٥٤}{ولا شعاع يؤنس ظلمته الباطنة؛ وغرفه مسكونة بشبح}{اسمه الظلام، ظل بارد أبكم، أقوى من جميع جند السماء.}{}

\Couplet{٥٥}{هذا الأنبوب، هذا المجرى الملغز الذي وُضعت نهايته قبل بدايته؛}{يطول ويتسع ويتقلص وينكسر؛ لغز وآلة وجهاز يتحرك من تلقاء نفسه.}{}

\Couplet{٥٦}{أول الأواني التي صنعها الفخّار بجانب موج كريسورواس الأزرق المخضر؛}{أحسبني أراه يبتسم إذ يرى أي جزاء قدّم للعالم!}{\BurtonNote{حاشية برتون: أبانا، نهر دمشق.}}

\Couplet{٥٧}{وكيف تكون الحياة باهتة، غير حقيقية، باطلة، كالمشاهد التي تترنح حول السكران؛}{وكيف لا يعني «الوجود» أن تكون، بل أن ترى وتسمع وتشم وتذوق وتحس.}{}

\Couplet{٥٨}{قطرة في مد المحيط الذي لا تحدّه حدود، في قفر من العذاب لا يُسبر؛}{حيث يعيش الملايين حياتهم المروعة بإماتة ملايين أخرى.}{}

\Couplet{٥٩}{وكيف يسعى الإنسان، بقلب يريد أن يصعد بالحب إلى الحب الكلي؛}{إلى اجتذاب فرصة جهنمية لتصيبه، كما تجتذب المآذن نار الصاعقة.}{}

\Couplet{٦٠}{وكيف يبني التراب على التراب برجًا وسورًا، لينهارا عند لمسة من الزمان؛}{وكيف يطمع التراب على التراب، من سهل شنعار، في تسلّق أعالي السماء.}{}

\Couplet{٦١}{ما أقصر هذه الحياة، وما أطولها مع ذلك؛ ما أكذب هنائها، وما أصدق ويلاتها؛}{نوبة حمى تتخللها تشنجات، تحدد مطلعها وختامها.}{}

\Couplet{٦٢}{آه، ما أبهج النهار بإشراق الشمس، وما أضوأ النسيم، وما أمرح الجموع}{التي التقت على ضفة النهر لتلعب، حين كنت شابًا، حين كنت شابًا.}{}

\Couplet{٦٣}{فرح كهذا يعم الجميع لا يمكن أن يبهت؛ ومع ذلك جاء الهمس البارد:}{وجه قد شحب، وجسد قد غاب؛ ترك الضفة وسبح عابرًا المجرى.}{}

\Couplet{٦٤}{وظل المحتفلون يرقصون ويغنون ويخطون على هذه الضفة من مد الزمان العميق؛}{وظلوا يغادرون، واحدًا واحدًا، ويسيرون إلى الضفة الأخرى البعيدة الضبابية.}{}

\Couplet{٦٥}{والآن أفلت آخرهم ليستكشف هناك قفر الموت الكئيب؛}{والآن لا يزال حاجّ واحد، منهك مهجور، يتلكأ على الشاطئ الوحيد.}{}

\Couplet{٦٦}{نعم، تقف الحياة في مد الشباب ساكنة؛ وفي الرجولة تجري برفق وبطء؛}{وانظر، حين تدنو من الغاية السحيقة، كيف تلمع المياه وتتدفق سريعة!}{}

\Couplet{٦٧}{والموت موتان: فالموت الذي نراه يتساقط كالورق في خريف عاصف؛}{أما موتنا، موتنا نحن، فعوالم منهارة، وكرة منهدمة، والنهاية الأخيرة لكل شيء.}{}

\Couplet{٦٨}{نعيش حياتنا مع محتالين وحمقى، أموات وأحياء، أحياء وأموات؛}{ونموت بين من يجس النبض ومن يضيق علينا الرأس ويكدّره.}{}

\Couplet{٦٩}{وآه، يا للأسف! ما نكاد نستوعب الدرس حتى يحين أجله القاتل؛}{فيأمرنا القدر أن نحزم كتبنا ونحملها، بما فيها وبأجسادنا، إلى الدود.}{}

\Couplet{٧٠}{ما نكاد نتعلم استعمال النصل حتى تشيخ المعصم وتتيبّس؛}{وما نكاد نتعلم إدارة القلم حتى يخور الفكر والخيال من البرد.}{}

\Couplet{٧١}{وما نكاد نجد طريق الحب، ونُغرق الذات، وننسى «أنا»؛}{حتى يقبض الشك الحزين على القلب، وحتى يبدأ الإنسان، الإنسان، في الموت.}{}

\Couplet{٧٢}{وما نكاد نتسلق ذرى الحكمة، ونبصر حولنا مشهد فسجة؛}{ونتنفس نسيم السماء، ونسمع نغم الأفلاك المتناسق؛}{}

\Couplet{٧٣}{حتى يجتاز راكب الجمل سريعًا الفلاة العاوية، وقد دفعه القدر؛}{وبتلويحة من عصاه السحرية يعجّل بالأحياء إلى اللحاق بالموتى.}{\BurtonNote{حاشية برتون: يمتطي الموت في بلاد العرب جملًا، لا حصانًا أشهب.}}

\Couplet{٧٤}{ما أشد العبء، وما أغرب الصراع؛ ما أكثر ما فيه من بهاء وعجب وخوف؛}{الحياة ذرة من ذلك الفضاء اللامتناهي الممتد بين هنا وهناك.}{}

\Couplet{٧٥}{وكيف يعجز الفكر عن استبصار السر الذي تحرسه الآلهة؛}{سر «لماذا» الميلاد والحياة والموت، ذلك الحجاب الإيزيسي الذي لا تستطيع يد أن تمزقه.}{}

\Couplet{٧٦}{تصنع الغدوات التي لا تنتهي يومنا؛ و«ما نحن عليه» دائمًا «ما سنكونه»، إلى أن}{يطبق الليل؛ إنه كله حلم، ومع ذلك نموت؛ ثم ماذا، ثم ماذا؟}{}

\Couplet{٧٧}{ولا يزال النسّاج يعمل بنوله، وسَداه ولُحمته الإنسان البائس؛}{ينسج الرسم المظلم الذي لا نموذج له، من فرط ظلمته نشك أن له تدبيرًا.}{}

\Couplet{٧٨}{ألا تخجل، أيها الصانع، حين تسمع، وسط عاصفة الدموع والدم؛}{الإنسان يقول إن رحمتك صنعت ما هو كائن، وإنك نظرت إلى المصنوع وقلت إنه حسن؟}{}

\Couplet{٧٩}{والعجب أن الإنسان يستطيع أن يبتسم وهو يحلم حلمه الشبحي المرعب؛}{لَخير منه تلك الذرة الغافلة التي تطن في شعاع الصباح!}{}

\Couplet{٨٠}{آه من لوعة حياتنا المروعة! كيف جرؤت، يا الله، أن تعبث هكذا}{بالحب والمودة والصداقة، بكل ما يكشف الإله في الطين الفاني؟}{}

\Couplet{٨١}{لكن آه، ماذا ينفع الإنسان أن ينوح؟ أتحقق الدموع ما لم تحققه الابتسامات؛}{وهل يولّد الاستغراق في الحزن خاطرًا من الفرح؟ آه، أخمد التنهد، وانسَ الخاطر!}{}

\Couplet{٨٢}{أسكت سعيك الذي لا أول له، واضبط شكوى طبيعتك العبثية؛}{لا أحد يصغي، ولا أحد يكترث لك أو لما يخصك؛ كم مثلك جاء ومضى؟}{}

\Couplet{٨٣}{كفّ، أيها الإنسان، عن النواح والبكاء والعويل؛ وتمتع بساعتك المضيئة في الشمس؛}{نرقص على حافة الموت الجليدية، ولكن أيقلّ بذلك ما في الرقص من مرح؟}{}
